\chapter{Introduction to Data Mining}
\label{ch:introduction-to-data-mining}

\section{Overview and Fundamental Concepts}

Data Mining represents the computational process of discovering patterns, anomalies, and actionable knowledge from vast repositories of data. It operates at the intersection of machine learning, statistics, and database systems.

\begin{definition}[Data Mining]
Data Mining is the non-trivial extraction of implicit, previously unknown, and potentially useful information from large volumes of structured or unstructured data.
\end{definition}

\section{The Knowledge Discovery (KDD) Process}

The standard Knowledge Discovery in Databases (KDD) pipeline consists of an iterative series of stages:
\begin{enumerate}
    \item \textbf{Data Cleaning:} Removing noise, outliers, and handling missing values.
    \item \textbf{Data Integration:} Combining data from multiple heterogeneous sources.
    \item \textbf{Data Selection and Transformation:} Selecting task-relevant attributes and applying aggregation or normalization.
    \item \textbf{Data Mining:} Applying intelligent algorithms to extract patterns and models.
    \item \textbf{Pattern Evaluation and Knowledge Presentation:} Interpreting patterns and translating findings into domain-specific decisions.
\end{enumerate}

\section{Core Data Mining Tasks}

Data mining tasks are commonly divided into two broad categories:
\begin{itemize}
    \item \textbf{Predictive Tasks:} Use existing variables to predict unknown or future values of target attributes (e.g., classification, regression).
    \item \textbf{Descriptive Tasks:} Find human-interpretable patterns describing the underlying data structure (e.g., clustering, association rule mining, anomaly detection).
\end{itemize}

\begin{example}[Association Analysis]
In a market basket analysis setting, the goal is to identify items frequently bought together. Given items $A$ and $B$, an association rule $A \implies B$ is evaluated through metrics such as support and confidence.
\end{example}
